\documentclass[10pt,a4paper]{amsart}
\usepackage[margin=19mm]{geometry}
\usepackage{amsmath,amssymb,mathtools}
\usepackage[scr=rsfs,cal=euler]{mathalfa}
\usepackage{xfrac}
\usepackage[unicode,hidelinks,bookmarks=false]{hyperref}
\newcommand{\ie}{i.e.\ }
\newcommand{\cf}{cf.\ }
\newcommand{\resp}{respectively\ }
\newcommand{\Subsection}{\subsection}
\providecommand{\nobreakdash}{\nobreak\hbox{-}}
\setlength{\emergencystretch}{2em}
\multlinegap0pt
\usepackage{xcolor}
\usepackage{mathtools}
\usepackage{amssymb}
\newtheorem{prop}{Propostion}
\newcommand{\littleo}[1]{\underset{#1}{o}}
\title{Rotation number for a homogeneous vector field in $\mathbf{C}^{2}$}
\author{Adrien Kachkachi}
\date{Source: arXiv:2605.26342v2 (CC BY 4.0)}
\begin{document}
\maketitle

\noindent\emph{Source and licence.} Rotation number for a homogeneous vector field in $\mathbf{C}^{2}$. Version arXiv:2605.26342v2 (CC BY 4.0), distributed by arXiv under the Creative Commons Attribution 4.0 International licence, http://creativecommons.org/licenses/by/4.0/, as recorded in the arXiv licence metadata for this exact version and shown on the abstract page https://arxiv.org/abs/2605.26342v2. Full licence notice: \texttt{licenses/arxiv-2605.26342-CC-BY-4.0.txt}.\par\smallskip

\noindent\textit{Editorial scope.} Counted unit: the article's prop prop:Dim\_Hausdorff\_modele (source0.tex lines 585--588). The article's complete original proof of it is delivered with it (source0.tex lines 614--647, one \texttt{proof} environment of 34 lines, which the article places away from the statement). No mathematics is written, repaired or re-derived: the statement and the proof are the article's own lines, in the article's own notation, and no retained line is substituted or re-flowed.  No further source-local context is needed: every cross-reference of every retained line resolves inside the two delivered spans.  Nothing was omitted from any retained span, so the package carries no omission record.  No retained line contains a citation, so the wrapper carries no bibliography and no bibliography entry.  No retained line contains a figure environment, an image-inclusion command or drawing code, so no image is delivered and the package declares no asset dependency.  Editorial formatting only, in addition: an AMS \texttt{amsart} preamble that restates the article's own declarations and macros used by the retained lines, namely the environments \texttt{prop} on separate counters, together with the standard packages those lines need.  The retained environments print in this document as Proposition 1 in order of appearance, whereas the article prints the same items with the numbering of its own full document.  The complete native source of the article as distributed in the arXiv e-print for this exact version is bundled unchanged as \texttt{sources/201410/source0.tex}.
\begin{prop}
\label{prop:Dim_Hausdorff_modele}
    The Hausdorff dimension of $\mathcal{K}$ is 0.
\end{prop}

\begin{proof}[Proof of Proposition \ref{prop:Dim_Hausdorff_modele}]
    A natural cover of $\mathcal{K}$ to consider is $\mathcal{K}_{n}:=\mathcal{I}\setminus\mathcal{H}_{n}$ for $n\in\mathbf{N}$. It is clear that $\mathcal{H}_{n+1}\setminus\mathcal{H}_{n}$ consists of $2^{n+1}$ disjoint intervals, and since $\mathcal{K}_{0}$ has two components, $\mathcal{K}_{n}$ is made of $2^{n+1}$ disjoint intervals. Let $w$ be a word of length $n$ and $\mathcal{J}$ one of the intervals intervals of $\mathcal{K}_{n}$. If $\delta_{j}^{w}$ denotes the proportion removed at the $j$-th step of the renormalization induced by $w$, the length of $\mathcal{J}$ is at most
    \[
    \prod_{j=0}^{n}(1-\delta_{j}^{w})\leq\frac{1}{2^{n+1}}.
    \]
    Now, let $l\in\mathbf{N}^{*}$, and set $\delta_{l}:=1-e^{-l^{2}}$. We consider the cover $\mathcal{K}_{lN}$ of $\mathcal{K}$ for some $N\in\mathbf{N}^{*}$, consisting of $2^{lN+1}$ intervals of length at most $1/2^{lN+1}$. There exists $n\in\mathbf{N}$ such that a finite word $w$ satisfies $\lvert H\left(w\right)\rvert/\lvert I\left(w\right)\rvert<\delta_{l}$ if and only if either the number of $L$ or the number of $R$ in $w$ is at most $n$. Therefore, if $lN\geq2n+1$, the number of words of length $lN+1$ satisfying $\lvert H\left(w\right)\rvert/\lvert I\left(w\right)\rvert<\delta_{l}$ is at most
    \[
    \sum_{j=0}^{n}\left(
    \begin{array}{c}
        lN+1 \\
        j 
    \end{array}\right)
    +\sum_{j=lN+1-n}^{lN+1}\left(
    \begin{array}{c}
        lN+1 \\
        j 
    \end{array}\right)
    =2\sum_{j=0}^{n}\left(
    \begin{array}{c}
        lN+1 \\
        j 
    \end{array}\right)=\frac{2}{n!}(lN)^{n}(1+\littleo{N\rightarrow+\infty}(1)).
    \]
    This implies that, among the $2^{(l+1)N+1}$ intervals composing $K_{(l+1)N}$, at least $2^{N}(2^{lN+1}-O(N^{n}))=2^{(l+1)N+1}(1-o(1))$ of them are smaller than $(1-\delta_{l})^{N}/2^{lN+1}$. This way, we obtain that, for any $\varepsilon>0$, the Hausdorff outer measure of dimension $\varepsilon$ bounded by $1/2^{(l+1)N+1}$ of $\mathcal{K}$ is less than
    \begin{equation}
    \label{Eq:Mes_Hausdorff}
      \frac{2^{1-\varepsilon}(lN)^{n}}{n!}(1+o(1))\left[2^{1-(l+1)\varepsilon}\right]^{N}+2^{1-\varepsilon}(1-o(1))\left[2^{l+1-l\varepsilon}(1-\delta_{l})^{\varepsilon}\right]^{N}.  
    \end{equation}
    For $\varepsilon$ larger than
    \[
    M_{l}:=\max\left(\frac{1}{l+1},\frac{(l+1)\log2}{l\log2-\log(1-\delta_{l})}\right),
    \]
    the quantity (\ref{Eq:Mes_Hausdorff}) goes to zero when $N$ tends to infinity, proving that $\text{dim}_{H}\mathcal{K}\leq M_{l}$ for all $l\in\mathbf{N}$. Finally, since $M_{l}\rightarrow0$ as $l\rightarrow+\infty$, we get $\text{dim}_{H}\mathcal{K}=0$.
\end{proof}

\end{document}
